%\title{iTi 2014 Turbulence Conference abstract}
\documentclass[11pt]{article}
%
%           Packages            %
%
\usepackage{amsmath}
% \usepackage[dvips]{epsfig}
\usepackage{graphicx}
\usepackage{geometry}
%
%           Page Settings           %
%
\pagestyle{empty} \geometry{left=2.5cm, right=3cm, nohead, nofoot,
top=2.5cm, bottom=1.5cm}
%
%                                   %
\begin{document}
%%%%%%%%%%%%%%%%%%%%%%%%%%%%%%%%%%%%%
%           Header                  %
%%%%%%%%%%%%%%%%%%%%%%%%%%%%%%%%%%%%%
% Abstract template for the         %
% iti Conference on Turbulence 2014 %
% October    21-24, 2014             %
% Bertinoro, Italy                  %
%%%%%%%%%%%%%%%%%%%%%%%%%%%%%%%%%%%%%
%
\begin{center}
    {\LARGE\bf Title}
\end{center}
%
\begin{center}
    \underline{Corresponding author} $^{1}$,
    B. Author  $^{2}$ and C.  Author $^{3}$\\
\end{center}
%
%           Authors                 %
%
{\small
$^{1}$ Department, University/Company, City, Country\\
$^{2}$ Department, University/Company, City, Country\\
$^{3}$ Department, University/Company, City, Country}\\
%
\vspace{0.2cm}\\
%
% Please specify two key-words e.g.:
key-words: {Fundamentals, Simulation}\\
%
\vspace{0.2cm}\\
%
%
%           Text                    %
%
{\bf Abstract:}\\
To obtain a standardized layout, we ask all authors to use the
given LaTex- or Word style file and to take notice of the
following instructions. The text should be written in English and
must not exceed one page. A second page is only allowed for figures.
Please, submit the PDF file of your abstract by using the provided form 
on the conference homepage 
\textit{http://www.fdy.tu-darmstadt.de/iti2014/iti2014\_1/Home.en.jsp}.
Deadline for submission of abstracts is April 20, 2014 .\\
\\
Examples:\\
These are links to equation \ref{eq_label} and to figure \ref{fig_label_for_the_figure}.\\
This is a link to reference \cite{Code_2}.\\
%
%           One Figure            %
%
\begin{figure}[h!]
    \begin{center}
        \includegraphics[height=1.5cm]{test_picture.eps}
        \caption{\emph{Caption for the figure}\label{fig_label_for_the_figure}}
    \end{center}
\end{figure}
%
%           Two Figures             %
%
\begin{figure}[th!]
    \begin{minipage}[h]{.49\linewidth}
        \begin{center}
            \includegraphics[height=1.5cm]{test_picture.eps}
            \caption{\emph{Caption for the figure on the left side}\label{fig_label_for_the_left_figure}}
        \end{center}
    \end{minipage}
    \hspace{.01\linewidth}
    \begin{minipage}[h]{.5\linewidth}
        \begin{center}
            \includegraphics[height=1.5cm]{test_picture.eps}
            \caption{\emph{Caption for the figure on the right side}\label{fig_label_for_the_right_figure}}
        \end{center}
    \end{minipage}
\end{figure}
%
%           Explanation             %
%
% "{.49\linewidth}"         the fraction of the paperwidth used for the left and right figure
% "\hspace{.01\linewidth}"  the fraction of the paperwidth used for the interspace between the
%                           two figures
% Hint:                     Instead of figures you can use the minipage also for equations or text.
%
%           Equation               %
\begin{align}
    a&=b+c\label{eq_label}\\
    d&=e+f \nonumber
    %%%%%%%%%%%%%%%%%%%%%
    \intertext{Text between equations}
    %%%%%%%%%%%%%%%%%%%%%
    x&=y+z
\end{align}
%
%           Explanation             %
%
% "&"               used for alignment for the different equations
% "\\"              used for linebreak
% "\nonumber"       used to cancel the numbering of this equation
% "\intertext"      used to place text between equations. The advantage is that you can
%                   align equations with text in between
%
\subsection*{Full paper submission:}
For the full paper submission the SVMult Class should be used, available on \\
\textit{http://www.fdy.tu-darmstadt.de/iti2014/submitabstractpaper/index.en.jsp}.

\begin{thebibliography}{1}
%
% Example for the citation of articles:
%
\bibitem{Code_1}
Author1 A., Author2 B., Author3 C. and Author4 D., (Year),
\newblock Title, \emph{Journal, vol. ..., pp. ...-....}\\
%
% Example for the citation of books:
%
\bibitem{Code_2}
Author1 A., Author2 B., Author3 C. and Author4 D., (Year):
\newblock \emph{Title}, Publisher
%
% % Example for the citation of papers in proceeding volumes:
%
\bibitem{Code_3}
Author1 A., Author2 B., Author3 C. and Author4 D., (Year):
\newblock Title, \emph{Proceedings}, City, pp. ...-...
%
\end{thebibliography}
%
\end{document}
%